\clearpage\thispagestyle{empty}
\pdfbookmark[0]{Quellen zur Juno-Anzeige vom Januar 1808}{juno1808elementsVII-sources}
\begin{center}\textbf{Quellen zur Juno-Anzeige vom Januar 1808}\end{center}
\noindent Diese KI-Übertragung ins heutige Deutsch folgt dem gesammelten Text \textit{Juno. Elemente VII.} in \href{https://gdz.sub.uni-goettingen.de/id/PPN236007467}{Gauß’ \textit{Werken}, Bd.~VI} (Göttingen, 1874), S.~292--293. Der Erstdruck erschien unter der Überschrift \textit{Göttingen} in den \textit{Göttingischen gelehrten Anzeigen}, Stück 14, 23. Januar 1808, \href{https://manifests.sub.uni-goettingen.de/iiif/image/gdz:PPN319721507_1808:00000131/full/1600,/0/default.jpg}{S.~129} und \href{https://manifests.sub.uni-goettingen.de/iiif/image/gdz:PPN319721507_1808:00000132/full/1600,/0/default.jpg}{S.~130}. Im Exemplar der Niedersächsischen Staats- und Universitätsbibliothek Göttingen steht auf S.~129 der Name Gauß am Rand. Die Anzeige spricht von Gauß in der dritten Person.
\par\medskip
Die Zeitschrift druckt die VII. Elemente in der Anzeige ab. Die Werkausgabe verweist in eckigen Klammern auf den gesonderten Brief an Bode vom 24. Januar 1808 auf S.~294. KI-Übertragung und KI-Satz.
\clearpage
